\chapter*{前言}
\addcontentsline{toc}{chapter}{前言}

线性代数是高等学校理工科各专业的一门重要基础课，它的研究对象是向量、矩阵、
线性方程组、线性变换与二次型等。线性代数不仅是学习后续课程（如概率统计、
数值分析、运筹学、机器学习等）的必备工具，其公理化与结构化的思想方法本身
也是数学思维训练的重要内容。

本书按照标准的教学大纲编写，全书共六章：

\begin{itemize}[itemsep=1pt]
  \item 第一章\quad 行列式——行列式的定义、性质、展开定理与克拉默法则；
  \item 第二章\quad 矩阵及其运算——矩阵运算、逆矩阵、分块矩阵、初等变换与矩阵的秩；
  \item 第三章\quad 向量组的线性相关性——线性组合、线性相关性、向量组的秩与向量空间；
  \item 第四章\quad 线性方程组——解的判定定理、基础解系与解的结构；
  \item 第五章\quad 特征值与特征向量——内积与正交性、方阵的对角化、实对称矩阵的正交对角化；
  \item 第六章\quad 二次型——二次型的标准形、惯性定理与正定二次型。
\end{itemize}

每章末配有习题，书末附录给出了部分习题的答案或提示。本书只假定读者具备高中
数学知识；第一章从二元、三元线性方程组出发引入行列式，各章内容循序渐进，
教师可根据学时适当取舍（例如第三章与第四章的顺序可以互换讲授）。

本先行版包含第一至第五章，并配有“问与答”框与几何直观插图，以帮助读者理解
概念的本质；第六章与习题答案附录将在后续版本中陆续补齐。

本书为草稿（v0.1），全部内容开放审阅。读者若发现概念表述、公式、例题演算或
习题答案中的任何问题，欢迎通过代码托管平台的问题跟踪（Issues）页面逐条指出，
并请注明章节与题号，以便修订。

\begin{flushright}
编者\quad 2026年10月
\end{flushright}
